% from ega4-2-fr.tex
\oldpage[IV-2]{91}

est dense dans $X_i$, il résulte de (4.1.1.3) que
$\dim(X_i)=\dim(U\cap X_i)$, donc on a
$\dim_x(X)=\sup_i(\dim(X_i))$. En outre, les idéaux premiers minimaux de
$\mathcal O_x$ correspondent aux points génériques des $X_i$, donc
$(0,16.1.1.1)$, on a
$\dim(\mathcal O_{X,x})=\sup_i(\dim(\mathcal O_{X_i,x}))$. On est donc ramené
au cas où $X$ est \emph{irréductible}; comme $X$ est biéquidimensionnel par
(5.2.1), on a
$\dim(X)=\dim(\overline{\{x\}})+\operatorname{codim}(\overline{\{x\}},X)$
$(0,14.3.5.1)$, et l'on sait que
$\dim(\overline{\{x\}})=\operatorname{deg.tr}_k k(x)$ par (5.2.1) et
$\operatorname{codim}(\overline{\{x\}},X)=\dim(\mathcal O_x)$ par (5.1.2).
\end{proof}

\begin{corollary}[5.2.4]
\label{IV.2.5.2.4-fr}
Soient $X$ un préschéma localement de type fini sur un corps $k$, $\mathcal L$
un $\mathcal O_X$-Module inversible, $f$ une section de $\mathcal L$ au-dessus
de $X$, telle que l'ensemble $Y$ des $x\in X$ tels que $f(x)=0$ soit rare
dans $X$. Alors on a $\dim(Y)\leq\dim(X)-1$; les deux membres de cette
inégalité sont égaux si $Y$ rencontre toute composante irréductible de
dimension maxima de $X$.
\end{corollary}

\begin{proof}
Si $(X_\alpha)$ est la famille des composantes irréductibles de $X$, on a
\[
\dim(Y)=\sup_\alpha(\dim(Y\cap X_\alpha))
\]
et l'on est donc ramené au cas où $X$ est irréductible ($Y\cap X_\alpha$
étant rare dans $X_\alpha$ puisque tout $X_\alpha$ a un intérieur non vide
dans $X$). On peut se borner au cas où $Y\ne\varnothing$; alors, pour tout
point maximal $x$ de $Y$, on a (puisque $X$ est biéquidimensionnel)
$\dim(\overline{\{x\}})=\dim(X)-\operatorname{codim}(\overline{\{x\}},X)$,
et puisque $Y$ est rare dans $X$, on a
$\operatorname{codim}(\overline{\{x\}},X)=1$ par (5.1.8); d'où le corollaire.
\end{proof}

\begin{remarks}[5.2.5]
\label{IV.2.5.2.5-fr}
\begin{enumerate}
\item[(i)] Contrairement à ce qui a lieu pour les $k$-préschémas algébriques,
un préschéma localement noethérien $X$ n'est pas nécessairement caténaire
(cf. (5.6.11)); il l'est cependant si tous ses anneaux locaux $\mathcal O_x$
sont des quotients d'anneaux réguliers $(0,16.5.12)$ et en particulier si $X$
est \emph{régulier} (I, 4.1.4). Toutefois, même un schéma (intègre)
$X=\operatorname{Spec}(A)$, où $A$ est un anneau \emph{régulier}, n'est pas
nécessairement \emph{biéquidimensionnel}, autrement dit $(0,14.3.3)$ les
codimensions dans $X$ des divers points fermés de $X$ ne sont pas
nécessairement les mêmes. Par exemple, soient $B$ un anneau de valuation
discrète, $\mathfrak m=B\pi$ son idéal maximal, $k=B/\mathfrak m$ son corps
résiduel, $K$ le corps des fractions de $B$; soit $A$ l'anneau de polynômes
$B[T]$. Il y a dans $A$ des idéaux maximaux \emph{de hauteur} $2$, par exemple
$(\pi)+(T)$; mais il y a aussi des idéaux maximaux \emph{de hauteur} $1$, par
exemple l'idéal principal $(\pi T-1)$: en effet, un idéal premier principal
$\ne0$ est de hauteur $1$ (5.1.8); d'autre part
$A/(\pi T-1)$ est isomorphe à l'anneau de fractions $B_\pi$
$(0_{\mathrm I},1.2.3)$, qui n'est autre ici que $K$, ce qui prouve que
l'idéal $(\pi T-1)$ est maximal et de hauteur $1$.

\item[(ii)] Lorsque $X$ est un préschéma localement de type fini sur un corps,
on a vu (4.1.1.3) que $\dim(U)=\dim(X)$ pour tout ouvert partout dense $U$
dans $X$. Ce résultat \emph{ne s'étend pas} aux préschémas noethériens
réguliers, même s'ils sont biéquidimensionnels. Par exemple, soient $B$ un
anneau de valuation discrète, $\mathfrak m$ son idéal maximal;
$X=\operatorname{Spec}(A)$ a deux points, l'idéal $(0)$ et $\mathfrak m$, ce
dernier étant le seul point fermé; on a $\dim(X)=1$, mais l'ensemble ouvert
$U=\{(0)\}$ est de dimension $0$ (cf. \S~10).
\end{enumerate}
\end{remarks}

\subsection{Dimension du support d'un Module et polynôme de Hilbert}
\label{IV.2.s5.3-fr}

